Daily stock returns alternate between calm and turbulent periods. Deep anomaly detectors
trained on raw returns therefore tend to flag whole turbulent periods rather than
individual abnormal days, and no deep, out-of-sample anomaly detection study exists for the
Dhaka Stock Exchange (DSE). This thesis proposes a volatility-aware approach. Returns are
standardized by the conditional volatility of GARCH, EGARCH or GJR models fitted on
training data only, and an LSTM autoencoder trained on the standardized returns is compared
with the same autoencoder trained on raw returns. Because the DSE has a short clean history,
detectors are trained on global market indices and transferred to the DSE, both without
retraining and after fine-tuning. Anomaly labels come from three sources fixed before any
model is trained: known market events, extreme standardized residuals and injected
anomalies. Detectors are evaluated with event-based F1 alongside point-adjusted F1 over five
random seeds, with time-based splits and a separate analysis of calm and volatile periods.
DSE data are taken from the DSEX index from 2013 and liquid stocks, with corrections for
floor-price periods and bonus-share adjustments. A preliminary run on S\&P~500 data from
1999 to 2018 supports the design: the volatility model with the best training fit failed
residual tests on the test period, while a rule based on GARCH-standardized residuals
flagged 17 days spread across the sample, against 54 days concentrated in crisis years for
a rule on raw returns.
